\documentclass{xStandalone}

\begin{document}
\begin{tikzpicture}[scale=1.5]

\foreach \y in {1,2,...,20}
{
    \draw[gray,ultra thin] (0,\fpeval{-13.6/(\y^2)}) -- ++(6,0);
}

\foreach \y/\t in {1/-13.60,2/-3.40,3/-1.51,4/-0.85,5/-0.54}
{
    \path (0,\fpeval{-13.6/(\y^2)}) node[left] {\y} ++(6,0) node[right] {$\t$};
}

\foreach \y/\x in {2/0.5,3/1.0,4/1.5,5/2.0}
{
    \draw[-latex,very thick] 
    (\x,\fpeval{-13.6/(\y^2)})
    --
    (\x,\fpeval{-13.6/(1^2)});
}

\foreach \y/\x/\t in {3/3.0/\alpha,4/3.5/\beta,5/4.0/\gamma}
{
    \draw[-latex,very thick] 
    (\x,\fpeval{-13.6/(\y^2)})
    node[above] {$\text{H}_{\t}$}
    --
    (\x,\fpeval{-13.6/(2^2)});
}

\foreach \y/\x in {4/5.0,5/5.5}
{
    \draw[-latex,very thick] 
    (\x,\fpeval{-13.6/(\y^2)})
    --
    (\x,\fpeval{-13.6/(3^2)});
}

\path (0,0) node[above right] {$n$} node[left] {$\infty$};
\path (6,0) node[above left] {$E/\si{eV}$} node[right] {$-0.00$};

\path 
(1.25,-13.6) node[below] {莱曼系}
(3.50,-3.40) node[below] {巴尔末系}
(5.25,-1.51) node[below] {帕邢系};

\end{tikzpicture}
\end{document}